\chapter{数据布局、局部性与预取}

数据布局决定程序怎样使用缓存、TLB 和内存带宽。很多高性能优化并不是改变算法复杂度，而是把数据摆成硬件更容易消费的形状。理解布局，是理解数组、结构体、队列、图和矩阵内核的共同基础。

\section{AoS 与 SoA}

Array of Structures 把一个对象的字段放在一起，适合经常同时访问完整对象的场景。Structure of Arrays 把同类字段连续存放，适合批量处理某个字段的场景。例如粒子模拟中，如果每轮只更新所有粒子的 x 坐标，SoA 更容易形成连续访问和 SIMD；如果每次处理一个粒子的全部属性，AoS 更直观。

选择布局不能只看代码好不好写，要看热路径读写哪些字段、是否需要向量化、是否有 false sharing、是否跨页跳转。工程中常见折中是 AoSoA，把数据按小块组织，既保留局部对象关系，又适合 SIMD 和缓存块。

\begin{lstlisting}[language=C++,caption={AoS 与 SoA 的访问差异}]
struct Particle {
    float x;
    float y;
    float z;
};

struct ParticleSoA {
    std::vector<float> x;
    std::vector<float> y;
    std::vector<float> z;
};
\end{lstlisting}

如果热路径只更新所有粒子的 \code{x}，SoA 版本会连续读取和写入 \code{x} 数组；AoS 版本每次还把 \code{y} 和 \code{z} 附近的数据带入 cache line。若热路径总是同时使用 x、y、z，AoS 的空间局部性反而更自然。

布局选择要从访问矩阵出发，而不是从结构体是否“面向对象”出发。可以把字段列成列，把热路径列成行，在每个格子里标出读、写、只读配置、冷路径调试。若一个热路径只读少数字段，SoA 或冷热分离更有优势；若多个字段总是一起使用，AoS 或 AoSoA 更清晰；若需要 SIMD，字段连续性和对齐会更重要；若需要并发写，线程到字段和槽位的映射必须避免 false sharing。

\begin{lstlisting}[numbers=none]
hot path: update position
  reads:  x, y, z, vx, vy, vz
  writes: x, y, z

hot path: compute bounding box
  reads:  x, y, z
  writes: thread local min and max

cold path: debug print
  reads: name, id, material, history
\end{lstlisting}

这张访问矩阵会自然推出布局：调试字段不应混在高频数值字段里；只在统计阶段使用的字段不应污染每帧更新；线程本地输出应分开存放，最后合并。

\section{预取}

硬件预取器擅长识别顺序访问和固定步长访问。连续数组扫描通常能被预取器提前加载，随机指针追踪则很难。软件预取可以在少数场景有效，但如果预取太早会挤出有用缓存，太晚则来不及隐藏延迟。

最好的预取优化通常不是手写预取指令，而是改变访问模式：把随机访问排序，按块处理，压缩节点，减少指针跳转，让硬件预取器能看懂程序。

软件预取只有在地址能提前算出、未来确实会访问、当前有足够独立工作、预取距离合适时才可能有效。指针追踪中，如果下一步地址必须等当前节点加载完成，预取很难发挥作用；批量处理多个游标时，才有机会提前为另一个游标发起预取。预取错误还可能污染缓存，把真正热的数据挤出去。

因此本书把预取看成最后一步，而不是第一步。先让数据连续，先分块，先减少对象分散，先测出内存延迟或带宽瓶颈，再考虑预取指令。很多工程代码的性能来自“让硬件预取器自然工作”，而不是手写大量 prefetch。

\section{冷热数据分离}

对象里不是所有字段都同样热。把热字段和冷字段混在一起，会让缓存带入很少使用的数据。冷热分离把热路径频繁访问的字段放在紧凑结构中，把调试信息、统计字段、名字、配置等冷数据放到旁边结构或间接存储。数据库执行器、游戏引擎和网络服务都常使用这种思路。

冷热分离也有代价：代码复杂度增加，对象关系不如原来直观。是否值得做，要看热路径字段访问比例和缓存压力。设计数据结构时，先写清哪个路径是热路径，哪个字段在热路径必读或必写。

\section{内存分配和对象布局}

小对象频繁分配会带来分配器开销、碎片、TLB 压力和缓存局部性问题。高性能系统常用 arena、对象池、批量分配或结构化生命周期来减少热路径分配。并发分配器还要避免全局锁和跨线程释放导致的远端缓存流量。

\begin{warningbox}
不要把“使用堆”简单等同于慢。真正的问题是生命周期是否清晰、分配是否在热路径、对象是否分散、是否跨线程释放、是否破坏局部性。
\end{warningbox}

对象池和 arena 的核心收益不是“分配函数少调用几次”这么简单。它们把生命周期相近的对象放在相近内存中，减少元数据开销和碎片，提高遍历局部性，也让释放可以批量发生。并发环境下，还可以按线程或 NUMA 节点维护本地池，减少全局分配器锁和远端释放。

代价同样要写清楚。Arena 适合批量释放，不适合单个对象精确析构；对象池可能保留大量空闲内存；跨线程归还对象可能重新引入同步；自定义分配器如果没有测试，可能比通用分配器更难维护。高质量教材不能只说“用内存池更快”，必须说明适用的生命周期形状。

\section{布局与并发}

多线程程序的数据布局还要考虑写共享。线程本地数据应尽量分开，跨线程只读数据可以共享，跨线程写数据要减少同一 cache line 上的冲突。并行归约、线程池队列和日志缓冲区都常常需要为每个线程准备独立槽位。

\section{结构体大小和 ABI}

结构体字段顺序影响 padding 和总大小。较大的结构体会降低缓存密度，也会增加拷贝成本。公共 ABI 中随意调整字段顺序可能破坏二进制兼容，因此高性能库常把内部热结构和外部稳定接口分开。内部结构按性能布局，外部接口保持稳定语义。

\section{布局重构流程}

数据布局优化应按流程进行。第一步，确认热路径和访问矩阵。第二步，估算每轮实际需要读写的字节数。第三步，检查对象是否连续、是否跨页、是否存在共享写。第四步，设计一个只改变布局、不改变算法语义的对照版本。第五步，用相同输入和相同绑定策略测量时间、cache、TLB 和扩展曲线。第六步，评估代码复杂度是否值得。

很多布局优化失败，是因为同时改了太多东西：换容器、改算法、改线程数、改内存分配、改输入规模。这样即使变快，也不知道原因。教材实验必须强调单变量对照，让读者学会建立因果证据。

\section{AoSoA：折中不是妥协}

AoS 和 SoA 常被讲成二选一，但真实系统经常使用 AoSoA。它把数据按小块分组，每个块内部采用 SoA，块与块之间保持对象批次关系。这样可以让一个块适配 SIMD 宽度或缓存块，同时又不会把整个系统改造成完全分离的字段数组。粒子系统、物理模拟和向量数据库里都能看到类似结构。

\begin{lstlisting}[language=C++,caption={AoSoA 的基本形态}]
constexpr std::int32_t PARTICLE_BLOCK_WIDTH = 8;

struct ParticleBlock {
    std::array<float, PARTICLE_BLOCK_WIDTH> x;
    std::array<float, PARTICLE_BLOCK_WIDTH> y;
    std::array<float, PARTICLE_BLOCK_WIDTH> z;
    std::array<float, PARTICLE_BLOCK_WIDTH> vx;
    std::array<float, PARTICLE_BLOCK_WIDTH> vy;
    std::array<float, PARTICLE_BLOCK_WIDTH> vz;
};

struct ParticleBlocks {
    std::vector<ParticleBlock> blocks;
    std::int32_t size = 0;
};
\end{lstlisting}

这里的块宽度不是魔法数字。它可能来自 SIMD lane 数、cache line 大小、对齐要求、尾部处理成本和代码复杂度。块太小，循环层次变多但复用不足；块太大，尾部浪费和缓存压力上升。AoSoA 的价值是提供一个调节旋钮，让布局能同时服务向量化和对象组织。

AoSoA 也有成本。随机访问第 \code{i} 个对象需要计算块号和块内偏移；插入删除比 \code{std::vector<Particle>} 更复杂；调试打印不如 AoS 直观；接口要隐藏布局细节，否则业务代码会到处写块索引。工程中通常把 AoSoA 封装在专用容器或视图后面，让热路径拿到高效访问，普通路径仍保持清晰 API。

\section{访问矩阵如何落到代码}

访问矩阵不是文档装饰，它应该直接影响类型设计。假设日志统计系统的每条记录有时间戳、用户 ID、事件类型、原始文本、解析状态和调试信息。热路径只需要用户 ID 和事件类型做聚合，原始文本只在错误报告时使用。如果把所有字段放在一个大结构里，聚合时每次都把冷字段附近的数据带入缓存。更好的设计是把热字段压成连续数组，把冷字段放在旁边，通过索引关联。

\begin{lstlisting}[language=C++,caption={冷热分离的记录批次}]
struct HotLogBatch {
    std::vector<std::uint64_t> user_ids;
    std::vector<std::int32_t> event_types;
};

struct ColdLogBatch {
    std::vector<std::string> raw_lines;
    std::vector<std::int32_t> parse_status;
};

struct LogBatch {
    HotLogBatch hot;
    ColdLogBatch cold;
};
\end{lstlisting}

这个设计牺牲了“一个对象包含所有字段”的直观性，换来热路径更少的数据移动。它还让并发更容易：聚合 worker 可以只读 \code{hot}，错误处理线程再查看 \code{cold}。若业务经常需要完整记录，冷热分离可能不值得；若绝大多数请求只走热路径，它就很自然。布局设计必须从访问频率和访问组合出发。

\section{索引优于指针的场景}

指针直接表达对象关系，但在大规模数据结构中，索引常常更适合性能和序列化。索引可以是 \code{std::uint32_t} 或 \code{std::uint64_t}，指向连续数组中的位置。它比指针更紧凑，容易持久化和跨进程传递，也让对象搬迁和压缩更容易。图的邻接表、编译器 IR、游戏 ECS 和列式存储都大量使用索引。

索引也有代价。每次访问需要多一步数组寻址；删除对象后要处理空洞、代际版本或稳定 ID；越界和悬空索引需要调试检查。若系统需要稳定引用，可以使用“索引加 generation”的句柄，避免旧索引误指向新对象。这个思想和无锁结构里的 ABA 类似：同一个位置再次被复用时，需要版本信息区分“看起来一样”和“语义上还是同一个对象”。

\begin{lstlisting}[language=C++,caption={索引句柄表达稳定引用}]
struct EntityId {
    std::uint32_t index = 0;
    std::uint32_t generation = 0;
};

struct EntitySlot {
    std::uint32_t generation = 0;
    bool alive = false;
};
\end{lstlisting}

这里使用明确位宽类型，是为了让内存布局和序列化边界清楚。系统教材中不能随手使用 \code{long long} 这类平台相关宽度。布局设计和数据类型选择是一件事：类型宽度会影响结构体大小、cache 密度、文件格式和网络协议。

\section{压缩、编码和计算成本}

有时减少内存流量比减少计算更重要。列式数据库和搜索系统常把整数压缩成变长编码、差分编码或位图。这样每个元素读取的字节数减少，但解码需要额外 CPU。是否划算取决于瓶颈：若系统受内存带宽限制，压缩可能提高吞吐；若系统受计算限制，压缩可能变慢。

这个权衡也出现在分布式 shuffle 中。压缩中间数据可以减少网络和磁盘 I/O，但会增加 CPU 时间和延迟。单机 cache 层面的“少搬字节，多做一点计算”和分布式层面的“少发网络，多做一点压缩”是同一个思想。教材应让读者从底层内存开始建立这种成本交换模型。

\section{并发写布局}

并发数据布局要先标出写者。字段是否共享，不看它是否在同一个结构体里，而看哪些线程写它、写入频率如何、是否落在同一 cache line 上。一个常见错误是把多个线程的统计字段放进同一个 \code{Stats} 数组，觉得每个线程写自己的下标就安全。结果每个字段很小，多个线程的字段落在同一条 cache line，形成 false sharing。

更好的设计是把写热字段按写者分开，并且让每个写者独占 cache line 或至少独占写热点。读者汇总时再扫描这些槽位。若读取频率很高，可以维护分层汇总；若读取只是指标展示，可以接受延迟。并发布局本质上是在写入吞吐、读取成本和内存占用之间取舍。

\begin{lstlisting}[language=C++,caption={线程本地统计槽位}]
struct alignas(64) WorkerStats {
    std::uint64_t tasks_executed = 0;
    std::uint64_t steal_attempts = 0;
    std::uint64_t empty_polls = 0;
};

std::vector<WorkerStats> stats(static_cast<std::size_t>(worker_count));
\end{lstlisting}

这类结构适合 worker 高频写自己的统计。它不适合每个请求都读取全局精确值。如果业务需要实时限流，统计槽位还要配合周期性汇总或原子发布。布局优化总是要回到语义。

\section{迁移旧代码的策略}

把成熟代码从 AoS 改成 SoA 或冷热分离，风险很高。直接大改会破坏接口、测试和调试习惯。稳妥策略是分阶段迁移。第一阶段写访问矩阵和性能证据，证明布局确实是瓶颈。第二阶段增加新布局的内部表示，但保留旧接口，通过 adapter 暴露兼容访问。第三阶段把热路径迁到新布局，保留 reference 对照。第四阶段逐步清理旧路径。

迁移时要防止“双写不一致”。如果旧结构和新结构同时存在，更新必须保证二者一致，或者明确一个是权威数据，另一个是缓存。缓存需要失效策略；双写需要事务边界；延迟同步需要版本号。布局重构不是单纯换容器，而是数据所有权迁移。

\section{标准容器的内存形状}

数据布局不只发生在自己定义的结构体里，也发生在每一次选择标准容器时。\code{std::vector} 的元素连续存放，适合顺序扫描、批量复制、SIMD 读取和按索引访问。它的扩容会搬迁元素，因此外部保存元素指针或引用时必须小心。\code{std::deque} 不是一个大连续数组，而是由多个固定大小块组成，头尾插入更稳定，但长距离顺序扫描的局部性通常不如 \code{std::vector}。\code{std::list} 每个节点独立分配，插入删除不移动其他节点，但遍历时会产生指针追踪、cache miss、TLB miss 和分配器开销。很多初学者只记住复杂度表，却没有把复杂度和硬件成本合起来看。

例如一个任务运行时维护 ready task。如果任务总数在构图后基本稳定，\code{std::vector} 加索引队列或环形缓冲通常更适合。如果任务需要频繁从中间删除，但遍历又很热，链表也未必正确，因为删除的理论优势可能被遍历成本吞掉。若真的需要稳定句柄，可以用索引加 generation，而不是直接把每个任务做成堆上节点。容器选择不是 API 风格问题，而是访问模式、生命周期和硬件局部性的共同结果。

\begin{lstlisting}[language=C++,caption={索引队列比链表更容易保持局部性}]
struct ReadyTask {
    std::uint32_t task_id = 0;
    std::uint32_t generation = 0;
};

class ReadyRing {
public:
    explicit ReadyRing(std::uint32_t capacity)
        : buffer_(capacity), head_(0), tail_(0), size_(0) {}

    bool push(ReadyTask task) {
        if (this->size_ == this->buffer_.size()) {
            return false;
        }
        this->buffer_[this->tail_] = task;
        this->tail_ = (this->tail_ + 1) % this->buffer_.size();
        ++this->size_;
        return true;
    }

    std::optional<ReadyTask> pop() {
        if (this->size_ == 0) {
            return std::nullopt;
        }
        const ReadyTask task = this->buffer_[this->head_];
        this->head_ = (this->head_ + 1) % this->buffer_.size();
        --this->size_;
        return task;
    }

private:
    std::vector<ReadyTask> buffer_;
    std::size_t head_;
    std::size_t tail_;
    std::size_t size_;
};
\end{lstlisting}

这个例子还不是并发队列，因为它没有锁和内存序。它要说明的是另一个层次：在考虑同步之前，先让数据本身连续、容量有界、访问路径可预测。后续给它加互斥或 SPSC 语义时，至少不会再被链式节点和频繁分配拖住。很多高性能结构的第一步不是“无锁”，而是“固定容量、连续存储、清楚所有权”。

\section{分配器不是只影响分配速度}

分配器常被理解成“申请内存快不快”，但它对布局的影响更深。通用分配器需要支持不同大小、不同生命周期、不同线程的分配和释放，因此会维护元数据、空闲链表、线程本地缓存和大块映射。小对象如果被零散分配，遍历时访问的是分配器历史留下的空间分布，而不是业务逻辑希望的顺序。一个图算法使用链式邻接节点，即使每个节点只有十几个字节，也可能因为节点分散而消耗大量 cache 和 TLB 资源。

Arena 的优势来自生命周期整齐。解析一个输入文件时产生的临时节点，可以在解析结束后一次性释放；构建一轮任务图时产生的任务对象，可以在图执行结束后批量销毁。这样做减少了单个对象释放的成本，也让相近生命周期的数据更容易靠近。缺点是对象不能随意提前释放，析构顺序要清楚，长期存活对象不能混进短生命周期 arena。否则 arena 会保留大量已经无用但不能回收的空间。

对象池则适合固定大小、频繁复用的对象，例如网络请求上下文、日志批次缓冲、任务描述符。对象池可以把对象初始化和内存申请分开，减少热路径分配；也可以按线程或 NUMA 节点分池，减少跨线程竞争。对象池的风险是复用对象时忘记清理旧状态，或者跨线程归还导致远端 cache line 流量。高质量实现要把 reset、owner thread、debug generation 和峰值容量都写进设计。

\begin{lstlisting}[language=C++,caption={固定大小对象池的基本接口形状}]
struct RequestContext {
    std::uint64_t request_id = 0;
    std::uint32_t attempt = 0;
    std::uint32_t worker_id = 0;
};

class RequestPool {
public:
    explicit RequestPool(std::size_t capacity) : storage_(capacity) {
        for (std::uint32_t index = 0;
             index < static_cast<std::uint32_t>(storage_.size());
             ++index) {
            this->free_.push_back(index);
        }
    }

    std::optional<std::uint32_t> acquire() {
        if (this->free_.empty()) {
            return std::nullopt;
        }
        const std::uint32_t index = this->free_.back();
        this->free_.pop_back();
        this->storage_[index] = RequestContext{};
        return index;
    }

    void release(std::uint32_t index) {
        this->storage_[index] = RequestContext{};
        this->free_.push_back(index);
    }

private:
    std::vector<RequestContext> storage_;
    std::vector<std::uint32_t> free_;
};
\end{lstlisting}

这个池子仍然只是单线程教学版本。并发版本必须决定 free list 由谁保护，是否每个 worker 一个本地池，跨线程释放是否进入 remote free 队列，是否要用 generation 防止旧句柄复用。不要因为代码短就忽略这些语义。内存池的工程难点不在“写一个 vector”，而在生命周期、并发所有权和调试可见性。

\section{面向数据的接口设计}

布局优化失败的常见原因，是内部换了 SoA，外部接口仍然要求每次返回一个完整对象。这样热路径又不得不临时拼对象，或者频繁跨数组取字段，收益被接口吃掉。面向数据的接口应该让调用者表达批量意图。例如不是提供 \code{get_particle(i)} 返回完整粒子，而是提供 \code{x_span()}、\code{velocity_span()} 或批处理函数，让热路径直接操作连续字段。

接口也要避免泄漏实现细节。完全暴露 \code{std::vector<float>&} 会让外部随意 resize，破坏布局不变量；完全隐藏成单元素 getter 又丢掉批量优化机会。折中做法是暴露只读或可写 \code{std::span}，并用类型说明生命周期和访问权限。只读输入、可写输出、临时 scratch、线程本地 partial 应该在函数签名中区分开。

\begin{lstlisting}[language=C++,caption={批量接口比单元素接口更能表达热路径}]
struct PositionView {
    std::span<float> x;
    std::span<float> y;
    std::span<float> z;
};

void integrate(PositionView position,
               PositionView velocity,
               float delta_time) {
    const std::size_t count = position.x.size();
    for (std::size_t i = 0; i < count; ++i) {
        position.x[i] += velocity.x[i] * delta_time;
        position.y[i] += velocity.y[i] * delta_time;
        position.z[i] += velocity.z[i] * delta_time;
    }
}
\end{lstlisting}

这个函数把布局意图表达给编译器和读者：三个坐标数组连续，循环边界统一，写入位置清楚。实际生产代码还要检查 span 长度一致、别名关系和对齐条件。接口设计的目标不是把所有优化都写进函数，而是不要一开始就把优化可能性封死。

\section{ECS 和列式布局}

Entity Component System 常见于游戏引擎，但它背后的思想适用于任何大量对象、少数字段热访问的系统。实体只是 ID，组件按类型连续存放，系统按组件批量处理。渲染系统关心位置和网格，物理系统关心位置和速度，AI 系统关心状态和目标。每个系统只扫描自己需要的组件，避免把完整对象的冷字段一起带入缓存。

日志分析和数据库执行器也使用类似思想。列式存储把同一列连续存放，过滤某个字段时只读取相关列；投影阶段再把需要的列组合输出。若查询只需要用户 ID 和事件类型，就不应该读取完整 JSON 文本。若后续需要原始文本，再通过行号或 offset 访问冷数据。列式布局的核心不是“数据库专属格式”，而是按访问组合移动最少字节。

列式布局还有压缩优势。同一列的数据类型相同、分布相近，更容易做字典编码、差分编码、位图和运行长度编码。压缩后减少内存和 I/O 流量，但增加解码成本。若过滤条件能直接在压缩表示上执行，收益很大；若每次都要完整解码，收益可能下降。布局、编码和算法必须一起设计。

\section{CSR：图结构的局部性案例}

图算法最容易暴露指针结构的问题。一个朴素邻接表如果每条边都是独立节点，遍历邻居时会不断指针跳转。Compressed Sparse Row 把所有边目标压到一个连续数组里，再用 offsets 数组标出每个顶点的邻居范围。这样遍历顶点 \code{v} 的邻居时，只需要扫描 \code{edges[offsets[v]]} 到 \code{edges[offsets[v + 1]]} 之间的连续区间。

\begin{lstlisting}[language=C++,caption={CSR 图的核心布局}]
struct CsrGraph {
    std::vector<std::uint32_t> offsets;
    std::vector<std::uint32_t> edges;
};

std::uint64_t degree_sum(const CsrGraph& graph) {
    std::uint64_t total = 0;
    for (std::size_t vertex = 0; vertex + 1 < graph.offsets.size(); ++vertex) {
        const std::uint32_t begin = graph.offsets[vertex];
        const std::uint32_t end = graph.offsets[vertex + 1];
        total += static_cast<std::uint64_t>(end - begin);
    }
    return total;
}
\end{lstlisting}

CSR 的代价也要讲清楚。它适合静态图或批量更新后的图，不适合高频单边插入删除。插入一条边可能需要移动后续大量边，或者维护增量结构再周期性重建。它还把顶点 ID 映射和边排序变成预处理问题。很多系统会使用双层结构：主图用 CSR 保持扫描性能，增量更新放在小的补丁区，定期合并。这个设计和 LSM-tree、列式数据库的 delta store 很像：热更新和冷扫描分层处理。

\section{块大小的选择方法}

分块是布局优化的核心技巧，但块大小不能拍脑袋。选择块大小时至少要考虑五个因素。第一，块内工作集是否能放进目标 cache 层级。第二，块内循环是否能形成足够 SIMD 和指令级并行。第三，块大小是否造成尾部处理过多。第四，块是否方便按线程或 NUMA 节点分配。第五，块元数据是否过多。

以矩阵乘法为例，一个块需要同时容纳 A 的子块、B 的子块和 C 的子块。若三者合起来超过 L1 或 L2，块内复用就会下降。以日志聚合为例，一个批次需要容纳 key、value、hash、输出缓冲和临时计数。批次太小，函数调用和调度成本高；批次太大，cache 复用差，背压粒度粗。块大小应通过候选集实验选择，例如从 1 KiB、4 KiB、16 KiB、64 KiB、256 KiB 开始，观察吞吐、cache miss、TLB miss 和尾延迟。

块大小还影响错误恢复。分布式计算中，任务块越大，调度开销越低，但失败重做成本越高；块越小，负载均衡更好，但调度和元数据成本更高。单机 cache block、线程池 task grain、分布式 partition chunk 其实是同一个问题在不同尺度上的版本。

\section{可测量的布局报告}

布局优化的报告至少应包含四类数据。第一，结构体大小和对齐：使用 \code{sizeof}、\code{alignof}，列出字段顺序和 padding。第二，访问字节数估算：每处理一个元素理论上读取和写入多少字节，哪些字节是冷字段。第三，硬件指标：时间、带宽、cache miss、TLB miss、分支和 IPC。第四，语义影响：接口是否改变，生命周期是否改变，是否影响稳定引用和并发写入。

只有耗时数字的报告不够。假设 SoA 比 AoS 快了 1.8 倍，你还要说明原因是否来自少读冷字段、SIMD 更好、预取更稳定、TLB 覆盖更高，还是因为新版本顺便少做了别的工作。若无法解释，就应该补实验。可以构造一个只访问全部字段的场景，如果 AoS 反而更快，说明 SoA 的优势确实来自字段子集访问，而不是神秘加速。

下面是一份布局报告的文字模板，但写报告时要填真实数据。

\begin{lstlisting}[numbers=none]
layout report:
  workload: aggregate event type from N records
  baseline: vector<Record>, hot fields mixed with raw line
  candidate: HotLogBatch plus ColdLogBatch
  semantic change: external record id remains stable
  bytes per hot record: baseline estimated 80, candidate 12
  metrics: time, cache misses, dTLB misses, bandwidth
  result: candidate improves hot aggregation, cold error path unchanged
\end{lstlisting}

\section{布局优化的反例}

布局优化也会失败。第一类失败是工作集很小，本来就全部在 cache 里，换布局只增加代码复杂度。第二类失败是热路径需要完整对象，SoA 反而让字段分散在多个数组中，增加地址计算和 cache line 数。第三类失败是业务频繁插入删除，CSR 或紧凑数组重建成本太高。第四类失败是并发访问模式改变，新布局把原来分散的写入集中到同一 cache line。第五类失败是接口没有同步迁移，内部布局高效，外部仍然单元素调用。

因此本章不是要求读者把所有对象都改成 SoA，而是要求读者学会问问题：热路径读哪些字段，写哪些字段，数据能否批量处理，生命周期是否成批，是否有稳定 ID 需求，是否多线程写，是否需要序列化，是否有恢复和调试要求。布局是系统设计，不是局部技巧。

\section{日志系统的布局案例}

贯穿项目的日志统计可以用一个具体布局案例串起来。原始日志行是冷数据，解析后的 event type、user id、timestamp 和 value 是热数据，解析错误信息是低频冷数据。若每条记录都保存成一个包含字符串和多个字段的大结构，聚合 event type 时会把大量原始文本和调试字段带入 cache。更好的设计是先把输入按 batch 解析，热字段放入列式数组，原始行或错误上下文按 offset 另存。

这种设计还方便错误处理。正常路径只扫描热字段；遇到解析错误时，把错误行号和原始 offset 写入冷错误列表；需要输出错误报告时再回到原始文本。热路径和冷路径分开后，性能和调试都更清楚。很多系统性能差，不是因为算法复杂，而是每次处理热字段都拖着一堆冷调试信息走。

\begin{lstlisting}[language=C++,caption={日志批次的热冷布局}]
struct ParsedLogHotBatch {
    std::vector<std::uint64_t> timestamps;
    std::vector<std::uint64_t> user_ids;
    std::vector<std::uint32_t> event_types;
    std::vector<std::int64_t> values;
};

struct ParsedLogColdBatch {
    std::vector<std::uint64_t> raw_offsets;
    std::vector<std::uint32_t> error_rows;
};
\end{lstlisting}

这个布局不要求原始输入立即复制。\code{raw_offsets} 可以指向 mmap 文件或输入缓冲中的位置，但生命周期要清楚：如果输入缓冲释放，offset 就不能再用于错误报告。布局和生命周期永远绑定在一起。

\section{ECS 变体：稀疏集合}

ECS 常用 sparse set 维护实体到组件位置的映射。Dense 数组保存实际组件，sparse 数组从 entity id 映射到 dense index。遍历组件时扫描 dense，局部性好；查询某个 entity 是否有组件时，用 sparse 快速定位。删除时可以用 swap-remove，把最后一个元素移到删除位置，再更新 sparse。这个结构牺牲稳定顺序，换取紧凑遍历。

如果业务需要稳定顺序，swap-remove 就不合适；如果 entity id 非常稀疏，sparse 数组可能很大；如果组件跨线程频繁增删，需要同步。再次说明，数据结构没有脱离语义的绝对优劣。Sparse set 适合大量实体、组件遍历热、增删可接受重排的场景。

\section{Columnar batch 和 vectorized execution}

数据库的 vectorized execution 通常以 batch 为单位处理列式数据。一个 batch 可能包含几千行，每个算子处理一个 batch 后交给下一个算子。Batch 太小，函数调用和调度成本高；batch 太大，cache 压力和延迟增加。Batch 模型是 row-at-a-time 和全量材料化之间的折中。

日志统计项目可以借鉴这个模型。输入解析成 batch，filter 处理 batch，histogram 处理 batch，shuffle 按 batch 发送。每个 batch 有热字段列、选择向量、错误列表和统计摘要。这样项目既能展示数据布局，也能展示 I/O 背压：下游慢时，上游 batch 队列满，读取暂停。

\section{布局和序列化的差异}

内存布局和磁盘/网络格式不必相同。内存中为了计算使用 SoA，网络上为了兼容和压缩使用长度前缀或 protobuf 类格式，磁盘上为了恢复使用稳定 manifest。直接把内存结构 dump 到磁盘，短期方便，长期会被 ABI、padding、字节序和版本拖垮。高质量系统会显式转换格式。

转换本身有成本，因此需要批处理和缓冲。一次转换一条记录，函数调用和分配成本高；一次转换一个 batch，可以顺序读取列、批量编码、批量校验。布局优化和 I/O 章节在这里相遇：好的内存布局也要服务序列化。

\section{缓存友好的哈希表}

哈希表是系统里常见的性能热点。链地址法每个 bucket 指向链表节点，容易指针追踪；开放寻址把元素放在连续数组中，探测时更 cache 友好，但删除、扩容和高负载因子要小心。Robin Hood hashing、SwissTable 类结构通过控制探测长度和元数据字节改善局部性。教材不要求实现这些成熟结构，但要让读者知道 \code{std::unordered_map} 不是唯一选择。

日志聚合如果使用 \code{std::unordered_map<std::string, std::int64_t>}，每个 key 可能分配字符串和节点，局部性较差。若先把 event type intern 成 \code{std::uint32_t}，再用数组或开放寻址表，热路径会更清楚。数据布局优化常常从“把复杂 key 转成紧凑 id”开始。

\section{内存布局的调试工具}

布局优化需要工具。可以用 \code{sizeof} 和 \code{alignof} 检查结构体大小；用编译器警告或工具查看 padding；用地址打印检查数组连续性；用 \cmd{perf stat} 观察 cache 和 TLB；用 heap profiler 看分配热点；用 sanitizer 检查越界和 use-after-free。布局优化如果没有调试工具，很容易变成猜测。

项目可以增加一个 layout report 命令，打印关键结构体大小、对齐、字段数量、batch 容量和估算字节数。这个命令不直接提升性能，但能让读者看到数据结构的物理形状。系统教材应让抽象类型落到字节。

\section{布局迁移的测试策略}

从 AoS 迁移到 SoA 时，测试要覆盖读写一致性。可以保留旧 AoS reference，随机生成记录，分别填充 AoS 和 SoA，执行同一聚合，比较结果。还要测试边界：空 batch、单元素、非对齐长度、错误行、超长字符串、冷热字段缺失。迁移期间如果保留双表示，要测试更新一个字段后另一个表示是否同步或明确无效。

性能测试则要比较至少三条路径：完整记录访问、只访问热字段、错误报告访问。SoA 可能在热字段聚合上快，但完整记录访问变慢。报告要承认这种取舍，而不是只展示最好场景。

\section{对齐分配器和实际收益}

有些内核希望输入按 32 字节或 64 字节对齐。可以使用对齐分配器或平台 API 分配内存，也可以让容器内部保证对齐。对齐能减少跨 cache line 访问和满足某些 SIMD 要求，但它不是万能优化。若访问本来连续且 CPU 对非对齐访问支持良好，收益可能很小。若切片从对齐数组中间开始，起始地址仍可能不对齐。

项目中可以把对齐作为实验参数。生成同样数据，分别使用普通 vector、对齐缓冲、故意偏移一个元素的视图。比较自动向量化和手写 chunked 版本。报告中写清对齐如何保证、如何验证、对尾部和偏移如何处理。

\section{Scratch buffer 管理}

高性能算法常需要 scratch buffer，例如 scan 的 block sums、partition 的 counts、top-k 的候选、解析的临时位置列表。每次函数调用都分配 scratch，会引入分配器成本和内存波动。更好的做法是由运行时或调用者提供 scratch，或者每个 worker 复用本地 scratch。

Scratch 复用要清楚生命周期和大小。一个 worker 处理过大 batch 后，scratch 可能扩容到很大，后续小任务继续占用内存。可以设置最大保留容量，超过后释放或归还池。Scratch buffer 是数据布局和资源管理的交叉点。

\section{Layout API 的封装}

内部使用 SoA 或 AoSoA，不代表外部调用者要知道所有数组。可以提供 batch view，让热路径拿 span，冷路径通过 record id 获取完整记录。封装的目标是同时保护布局不变量和提供批量访问能力。单元素 getter 方便但可能慢；完全暴露 vector 容易破坏不变量。View 是折中。

\begin{lstlisting}[language=C++,caption={批量视图封装布局}]
struct EventColumnsView {
    std::span<const std::uint32_t> event_types;
    std::span<const std::int64_t> values;
};

EventColumnsView event_columns(const ParsedLogHotBatch& batch) {
    return EventColumnsView{batch.event_types, batch.values};
}
\end{lstlisting}

这个接口让聚合内核只看到自己需要的列，不依赖完整 batch 结构。后续如果 batch 从 vector 改成 mmap 或压缩列，接口仍有机会保持稳定。

\section{布局章节总结}

本章的主线是让数据形状服务访问形状。顺序热访问需要连续，字段子集热访问需要列式或冷热分离，高频并发写需要分片和对齐，稳定引用需要索引和 generation，静态图需要 CSR，动态更新需要增量结构。布局优化不是把所有东西改成一种格式，而是让每条热路径少搬无用字节、少追随机指针、少共享写。

设计布局时，先写访问矩阵，再写生命周期，再写并发写者，再写序列化边界，最后才写代码。这个顺序能避免“为了性能换结构，最后语义乱了”的常见错误。

\section{布局决策表}

选择布局时，可以先填表。字段是否总是一起访问？若是，AoS 可能清晰；若只访问少数字段，SoA 或冷热分离更好。对象是否需要稳定地址？若需要，考虑索引句柄或对象池；若不需要，连续数组更简单。是否高频插入删除？若是，CSR 这类紧凑静态结构可能不合适。是否要 SIMD？字段连续和对齐更重要。是否要序列化？内存布局和 wire format 要分开。是否多线程写？写者要分片或对齐。

这张表能避免布局争论变成风格争论。面向对象、面向数据、列式、池化都只是工具。真正决定工具的是访问矩阵、生命周期和并发语义。

\section{案例：任务图的布局}

任务图也需要布局。最直观的表示是每个节点一个对象，每个对象有 \code{std::vector} successors。小图这样写清楚，大图会有很多小 vector 分配，遍历后继时指针跳转多。CSR 式任务图把节点和边分别放入连续数组，ready count 在节点数组中，successor id 在边数组中。执行时，worker 完成节点后顺序扫描后继范围，局部性更好。

缺点是动态修改困难。如果任务图在执行中不断新增边，CSR 重建成本高。可以使用两层结构：静态主图用 CSR，动态新增任务用小的 pending edge 列表，stage 结束后合并。这个设计和静态图加增量边、列式主存加 delta store 是同一模式。布局思维可以迁移到运行时元数据。

\section{案例：模型权重的预告}

虽然 AI 在第三册展开，但模型权重布局可以作为预告。权重通常是大块只读数据，推理时被反复扫描；量化权重可能按 block 存放 scale、zero point 和 int8 数据；算子希望连续读取并适配 SIMD。这个场景和本章的列式、AoSoA、对齐、冷热分离完全相关。第二册学习日志 batch 和矩阵 block，不是只为日志系统，而是为后续权重和 tensor 布局打基础。

这也说明布局知识具有迁移性。字段从 event type 换成 tensor value，原则不变：热数据连续，元数据分离，block 大小匹配 cache 和 SIMD，序列化格式与内存计算格式分开。

\section{本章落到贯穿项目}

Compute Systems Engine 可以在本章加入两个实验模块。第一个是记录批处理模块：用 AoS、SoA 和冷热分离三种布局统计 key，比较热字段扫描、错误记录访问和内存占用。第二个是 worker stats 模块：紧凑统计槽位与 \code{alignas(64)} 槽位对照，观察多线程写入扩展性。这两个模块分别训练“读局部性”和“写共享”两种布局能力。

项目代码不必一口气实现复杂 AoSoA 容器，但要把布局实验和 benchmark harness 接起来。每个布局版本必须使用同一份输入、同一个 reference 结果和同一组线程策略。否则布局差异会和算法差异混在一起。

\section{从业务记录推导列式批次}

很多布局优化失败，是因为一开始就讨论“要不要 SoA”，却没有先讨论业务记录到底怎样被使用。正确推导应从一个具体记录开始。假设日志记录包含时间戳、用户编号、事件类型、数值、原始行、解析错误、调试标签和来源文件。解析阶段需要原始行、来源文件和错误字段；聚合阶段只需要用户编号、事件类型和数值；错误报告阶段需要原始行和错误字段；调试查询阶段才需要标签。把这些字段全部放在一个大对象里，聚合阶段每处理一条记录都会把大量冷字段一起带进缓存。聚合代码看起来只读三个字段，硬件实际搬动的字节却远多于三个字段。

推导列式批次时，可以先保留最直观的 AoS reference。它的价值是语义清晰，便于测试新布局。然后为热路径建立列式批次：用户编号是一列，事件类型是一列，数值是一列，解析状态可以是一列，原始行和调试信息放入冷批次。热聚合只接收热批次视图，错误报告才接收冷批次视图。这样做并不是为了追求某个抽象名称，而是为了让热路径少搬无用字节，让编译器更容易生成顺序扫描，让硬件预取器更容易工作。

迁移时要特别注意记录编号。列式批次里第 i 个用户编号、第 i 个事件类型、第 i 个数值和冷批次第 i 个原始行必须表示同一条逻辑记录。这个关系可以通过统一长度、统一追加、统一过滤掩码和稳定 record id 保持。如果某个阶段删除错误记录，不能只从一列删除而忘记其他列。更稳的做法是先生成 selection vector，表示哪些记录有效，后续聚合按 selection vector 读取；等到批次稳定后，再做 compact。这样能避免列之间长度和顺序不一致。

列式批次还要考虑异常和边界。解析一行失败时，热列是否填默认值？错误列如何记录原因？聚合阶段是否跳过失败记录？如果错误记录很多，selection vector 会不会成为新的热点？如果原始行很长，冷批次是否保存完整字符串，还是保存文件偏移和长度？这些问题都属于布局设计，而不是外围细节。真正的工程布局必须同时说明成功路径、失败路径、调试路径和序列化路径。

\section{稀疏结构和图布局}

图和稀疏矩阵是布局思想最集中的场景。最直观的图表示是每个节点一个对象，每个对象保存一个后继列表。这种写法适合教学和小图，但在大图上会产生大量小分配。遍历一个节点的邻居时，程序先读节点对象，再读 vector 元数据，再跳到邻居数组；如果每个邻居数组都来自不同分配位置，cache 和 TLB 压力会很高。CSR 表示把所有边放到一条连续边数组里，再用 offset 数组记录每个节点的边范围。遍历节点 v 的邻居时，只需要顺序扫描 edges 从 offsets[v] 到 offsets[v + 1] 的区间。

CSR 的核心收益是把指针追踪变成连续区间扫描。连续扫描让硬件预取器能提前加载，也让多线程分片更容易。若节点范围按编号切分，每个 worker 扫描一段节点和对应边，读路径相对稳定。代价是动态更新困难。插入一条边可能需要移动后续大量边，删除也可能留下空洞。因此 CSR 更适合静态图、阶段性重建图或批处理图。动态图可以使用主 CSR 加 delta 边列表：主图负责绝大多数只读遍历，新增边暂存在小结构中，阶段边界再合并。

稀疏矩阵也类似。按行压缩适合行遍历和矩阵向量乘；按列压缩适合列遍历；块稀疏适合局部密集区域和 SIMD。选择格式时，不能只说“稀疏矩阵用 CSR”。要看内核按行访问还是按列访问，是否需要转置，是否需要并行写输出，非零分布是否倾斜，行长度差异是否很大。行长度极不均匀时，按行分配线程可能负载不均；可以按非零数量切分，也可以把长行拆成多个任务，但拆长行又会引入输出累加同步。布局、算法和并行调度在这里不能分开。

Compute Systems Engine 后续做 shuffle、join 和任务图时，都可以借用稀疏结构思维。任务图的 successors 数组就是边数组，任务的 successor range 就是 offset。分区后的 key 到记录位置，也可以看成倒排索引或稀疏结构。读者如果只把 CSR 当成图算法专用格式，就错过了更重要的思想：当大量对象之间的关系可以静态收集时，把关系压成连续区间，通常比每个对象挂一串指针更适合现代硬件。

\section{布局实验报告怎样写}

布局实验报告必须说明输入规模、记录大小、字段访问比例、线程策略、内存分配方式和测量指标。只报告“快了两倍”没有意义，因为读者不知道快在哪里。一个合格报告至少要写三组结果：完整对象访问、热字段访问、冷路径访问。SoA 可能让热字段聚合更快，但完整对象调试变慢；冷热分离可能减少缓存流量，但增加索引和封装复杂度。报告要把收益和代价都列出来。

指标上，时间只是起点。还应观察 cache miss、TLB miss、内存带宽、分配次数、峰值内存、输出是否一致。多线程布局实验还要看扩展曲线和 false sharing。若紧凑统计槽位在单线程下更快，对齐槽位在多线程下更快，这不是矛盾，而是写共享形状改变了瓶颈。单线程结论不能直接推到多线程，多线程结论也不能忽略内存占用。

报告还要防止把算法变化伪装成布局变化。AoS 版本如果使用普通循环，SoA 版本同时换成 SIMD 和分块，就无法判断收益来自哪里。更严谨的方式是逐步改变：先只改布局，算法保持一致；再加分块；再加 SIMD；再加线程。每一步都有 reference 结果和性能指标。这样得到的不是一组漂亮数字，而是一条可复查的因果链。

\section{本章习题}

\begin{exercisebox}
习题一：为一个日志统计系统写访问矩阵。字段至少包括时间戳、用户 ID、事件类型、原始文本、解析错误、调试标签。热路径至少包括解析、按用户聚合、错误输出和调试查询。根据矩阵给出布局方案。

习题二：把一个指针链表设计改成索引数组设计。说明索引宽度、generation、删除策略、遍历局部性和序列化影响。不要只写“索引更快”，要写清生命周期如何保证。

习题三：设计一个冷热分离迁移计划。要求保留旧接口、建立 reference 测试、说明双写或缓存一致性策略，并给出何时删除旧布局的条件。

习题四：实现或设计 worker stats 的紧凑版和对齐版。报告每个版本的结构体大小、对齐、cache line 分布、写入线程和汇总成本。
\end{exercisebox}

\begin{labbox}
实验：实现 AoS 和 SoA 两个版本的字段求和。比较顺序访问和随机访问。再构造两个线程分别写相邻字段和分离字段的场景，观察 false sharing。
\end{labbox}
